%==============================================================================
% Chapitre 5 - Mécanique des fluides appliquée
%==============================================================================

\section{Introduction}

La mécanique des fluides est la branche de la physique qui étudie le
comportement des liquides et des gaz.

Dans le domaine de la balistique, le fluide qui nous intéresse
principalement est l'air.

Après sa sortie du canon, un projectile doit traverser l'atmosphère.
Son mouvement n'est donc pas uniquement influencé par la gravité.

L'air exerce également des forces qui modifient sa vitesse et sa
trajectoire.

Ces effets deviennent rapidement prépondérants dès que la distance de tir
augmente.

\begin{aretenir}

Dans le vide, un projectile est principalement soumis à la gravité.

Dans l'atmosphère, il est également soumis à l'action de l'air.

\end{aretenir}

\section{L'air est un fluide}

Dans la vie quotidienne, l'air semble souvent immatériel.

Pourtant, il possède :

\begin{itemize}
\item une masse ;
\item une densité ;
\item une pression ;
\item une température.
\end{itemize}

Il peut être comprimé, déplacé et accéléré.

L'air est donc un fluide au même titre que l'eau, même si ses propriétés
sont différentes.

Lorsqu'un projectile traverse l'atmosphère, il doit déplacer une certaine
quantité d'air devant lui.

Cette interaction est à l'origine de nombreux phénomènes balistiques.

\section{Pression}

La pression représente la force exercée par un fluide sur une surface.

Elle s'exprime dans le système international en pascals (Pa).

Au niveau moyen de la mer, la pression atmosphérique standard vaut
environ :

\begin{equation}
P \approx 101,325~\mathrm{Pa}
\end{equation}

Cette pression diminue avec l'altitude.

\begin{exemple}

Au sommet d'une montagne, la pression atmosphérique est plus faible qu'au
niveau de la mer.

L'air est donc moins dense.

Un projectile y subira généralement moins de traînée.

\end{exemple}

\section{Densité}

La densité de l'air joue un rôle majeur en balistique.

Plus l'air est dense :

\begin{itemize}
\item plus il freine le projectile ;
\item plus les effets aérodynamiques sont importants ;
\item plus la portée utile diminue.
\end{itemize}

Inversement, un air moins dense oppose moins de résistance au mouvement.

La densité dépend principalement :

\begin{itemize}
\item de la pression ;
\item de la température ;
\item de l'humidité.
\end{itemize}

\begin{aretenir}

Deux tirs identiques peuvent produire des résultats différents si les
conditions atmosphériques changent.

\end{aretenir}

\section{Température}

La température influence directement la densité de l'air.

Lorsque l'air se réchauffe :

\begin{itemize}
\item il se dilate ;
\item sa densité diminue.
\end{itemize}

Un air chaud est donc généralement moins dense qu'un air froid.

Pour un projectile, cela signifie :

\begin{itemize}
\item moins de traînée ;
\item une vitesse mieux conservée ;
\item une trajectoire plus tendue.
\end{itemize}

\begin{exemple}

Une munition tirée à 35°C et la même munition tirée à 0°C ne produiront
pas exactement la même trajectoire.

\end{exemple}

\section{Humidité}

L'influence de l'humidité est souvent surestimée.

Contrairement à une idée répandue, un air humide n'est pas nécessairement
plus dense qu'un air sec.

Dans les conditions atmosphériques courantes, l'effet de l'humidité sur la
trajectoire reste généralement faible par rapport :

\begin{itemize}
\item à la température ;
\item à la pression ;
\item au vent.
\end{itemize}

\section{Vitesse du son}

Le son se propage dans l'air à une vitesse finie.

Dans les conditions atmosphériques standard, cette vitesse vaut
approximativement :

\begin{equation}
c \approx 343~\mathrm{m/s}
\end{equation}

Cette valeur varie notamment avec la température.

La comparaison entre la vitesse d'un objet et la vitesse du son conduit à
la notion de nombre de Mach.

Cette notion jouera un rôle essentiel dans l'étude des projectiles
supersoniques.

\begin{rigueur}

La vitesse du son dans l'air sec peut être approximée par :

\begin{equation}
c \approx 331 + 0,6\,T
\end{equation}

où :

\begin{itemize}
    \item $c$ est exprimée en m/s ;
    \item $T$ est la température en °C.
\end{itemize}

\end{rigueur}

\section{Écoulements}

Lorsqu'un projectile se déplace dans l'air, il perturbe l'écoulement du
fluide autour de lui.

Selon sa vitesse, différents régimes peuvent apparaître :

\begin{itemize}
\item écoulement subsonique ;
\item écoulement transsonique ;
\item écoulement supersonique.
\end{itemize}

\begin{exemple}

\begin{itemize}
    \item Subsonique : $M < 0,8$
    \item Transsonique : $0,8 \le M \le 1,2$
    \item Supersonique : $M > 1,2$
\end{itemize}

\end{exemple}

Chaque régime possède ses propres caractéristiques aérodynamiques.

Les projectiles militaires modernes évoluent généralement dans le domaine
supersonique pendant une grande partie de leur vol.

\section{La résistance de l'air}

Toute personne ayant déjà sorti la main par la fenêtre d'un véhicule en
mouvement a ressenti la résistance de l'air.

Cette force s'oppose au mouvement.

Le même phénomène agit sur un projectile.

Cette résistance dépend notamment :

\begin{itemize}
\item de sa vitesse ;
\item de sa forme ;
\item de sa surface frontale ;
\item de la densité de l'air.
\end{itemize}

À mesure que la vitesse augmente, cette résistance devient rapidement très
importante.

\section{Notion de traînée}

La traînée est la force aérodynamique qui s'oppose au déplacement d'un
objet dans un fluide.

Elle constitue l'une des forces les plus importantes en balistique
extérieure.

Sans traînée :

\begin{itemize}
\item les projectiles conserveraient leur vitesse beaucoup plus
longtemps ;
\item les portées seraient nettement supérieures ;
\item les calculs balistiques seraient considérablement simplifiés.
\end{itemize}

L'étude détaillée de la traînée fera l'objet d'un chapitre spécifique.

\section{Pourquoi la forme du projectile est importante}

Tous les projectiles ne traversent pas l'air avec la même efficacité.

Leur géométrie influence fortement :

\begin{itemize}
\item la traînée ;
\item la stabilité ;
\item la conservation de la vitesse ;
\item la portée.
\end{itemize}

C'est pourquoi les formes modernes diffèrent fortement des anciennes
balles sphériques.

\begin{simulation}

Dans un moteur balistique simplifié, l'atmosphère peut être représentée
par une densité constante.

Les modèles plus avancés prennent en compte :

\begin{itemize}
\item la température ;
\item la pression ;
\item l'humidité ;
\item les variations avec l'altitude.
\end{itemize}

Le niveau de réalisme dépend directement de la précision souhaitée.

\end{simulation}

\section{Application à la balistique}

La mécanique des fluides intervient dans :

\begin{itemize}
\item la perte de vitesse ;
\item la portée ;
\item la dérive due au vent ;
\item les phénomènes transsoniques ;
\item la stabilité du projectile.
\end{itemize}

Elle constitue l'un des fondements de la balistique extérieure moderne.

\section{À retenir}

\begin{aretenir}

\begin{itemize}
\item L'air est un fluide possédant une masse et une densité.
\item La densité de l'air dépend notamment de la pression et de la
température.
\item L'air exerce une résistance au mouvement des projectiles.
\item Cette résistance est appelée traînée.
\item La forme du projectile influence fortement son comportement
aérodynamique.
\end{itemize}

\end{aretenir}

